\documentclass[11pt,a4paper]{article}

\usepackage[utf8]{inputenc}
\usepackage[T1]{fontenc}
\usepackage{amsmath,amssymb,amsthm}
\usepackage{graphicx}
\usepackage{booktabs}
\usepackage{geometry}
\usepackage{listings}
\usepackage{xcolor}
\usepackage[hidelinks]{hyperref}
\geometry{margin=2.5cm}

\lstset{
  basicstyle=\ttfamily\small,
  breaklines=true,
  frame=single,
  backgroundcolor=\color{gray!8},
  columns=fullflexible,
}

\title{\textbf{simgen}: A Prompt-Driven Framework for Generating\\
       Mathematical and Physical Simulations,\\
       Visualizations, and Reports}
\author{Generated and maintained with \texttt{simgen}}
\date{\today}

\begin{document}
\maketitle

\begin{abstract}
We present \texttt{simgen}, a software framework that converts a short
natural-language description of a mathematical or physical phenomenon into a
complete, runnable mini-project: a numerical simulation, publication-quality
plots, a \textsc{Manim} animation, and a compilable \LaTeX{} report. The system
uses a large language model (the Claude API) to translate a prompt into a strict,
validated specification, which a deterministic code generator renders into
portable artifacts built upon a reusable numerical core. A built-in library of
five validated reference systems (the Lorenz system, the double pendulum, the
simple harmonic oscillator, the Duffing oscillator, and the Van der Pol
oscillator) provides an offline fallback and a deterministic test bed. We
describe the architecture, the numerical methods, the code-generation pipeline,
and the validation strategy, and we demonstrate quantitative correctness through
analytic comparisons, energy-conservation checks, and a Lyapunov-exponent
estimate that recovers the textbook value for the Lorenz attractor.
\end{abstract}

\tableofcontents

\section{Introduction}
Educational simulation projects traditionally focus on a single phenomenon.
\texttt{simgen} instead automates the \emph{construction} of such projects: the
user provides a prompt, and the framework emits the code, figures, animation,
and report. The design separates three concerns: (i) \emph{interpretation} of
the prompt into a structured specification; (ii) \emph{deterministic rendering}
of that specification into artifacts; and (iii) a \emph{reusable scientific
core} of integrators and analysis tools. This separation makes the generated
output reproducible and testable, and it allows the whole pipeline to run
offline against a reference library when no model is available.

\section{Historical Context}
The numerical solution of ordinary differential equations has a long history,
from Runge's 1895 method~\cite{runge1895} to the embedded Runge--Kutta pairs of
Dormand and Prince~\cite{dormand1980}. The qualitative theory of dynamical
systems was transformed by Lorenz's 1963 discovery of deterministic
chaos~\cite{lorenz1963}, by the geometric viewpoint of Ruelle and
Takens~\cite{ruelle1971}, and by the universality results of
Feigenbaum~\cite{feigenbaum1978}. Practical diagnostics for chaos --- notably
Lyapunov-exponent estimation --- were developed by Benettin et
al.~\cite{benettin1980}, Wolf et al.~\cite{wolf1985}, and Rosenstein et
al.~\cite{rosenstein1993}. The reference systems shipped with \texttt{simgen}
are drawn from this canon~\cite{strogatz2015,guckenheimer1983}.

\section{Theory}
Every phenomenon is modelled as a first-order system of ordinary differential
equations,
\begin{equation}
  \frac{d\mathbf{y}}{dt} = \mathbf{f}(t, \mathbf{y}), \qquad
  \mathbf{y}(t_0) = \mathbf{y}_0, \quad \mathbf{y} \in \mathbb{R}^n .
\end{equation}
Higher-order equations are reduced to this form by introducing auxiliary
(velocity) variables, and spatial partial differential equations are handled by
the method of lines~\cite{leveque2007}, which discretises space to yield a large
ODE system. This uniform representation lets a single integrator and a single
code template serve an open-ended family of problems.

\section{Mathematical Derivations}
\subsection{Reduction of a second-order system}
A mechanical system $m\ddot{x} = F(t, x, \dot{x})$ becomes
\begin{equation}
  \dot{x} = v, \qquad \dot{v} = \frac{1}{m} F(t, x, v),
\end{equation}
with state $\mathbf{y} = (x, v)$. The simple harmonic oscillator is the linear
case $F = -m\omega^2 x$, giving the conserved energy
$E = \tfrac{1}{2}v^2 + \tfrac{1}{2}\omega^2 x^2$.

\subsection{The Lorenz system}
The Lorenz equations~\cite{lorenz1963},
\begin{align}
  \dot{x} &= \sigma(y - x), \\
  \dot{y} &= x(\rho - z) - y, \\
  \dot{z} &= xy - \beta z,
\end{align}
have fixed points at the origin and, for $\rho > 1$, at
$C_\pm = (\pm\sqrt{\beta(\rho-1)}, \pm\sqrt{\beta(\rho-1)}, \rho-1)$. For the
classical parameters $\sigma = 10$, $\beta = 8/3$, $\rho = 28$ all fixed points
are unstable and trajectories approach a strange attractor.

\section{Numerical Methods}
\texttt{simgen} provides fixed-step explicit Euler and classical fourth-order
Runge--Kutta (RK4) integrators, together with a bridge to SciPy's adaptive
solvers (RK45, DOP853, Radau, BDF, LSODA)~\cite{virtanen2020scipy,hairer1993}.
RK4 advances the solution by
\begin{equation}
  \mathbf{y}_{n+1} = \mathbf{y}_n + \tfrac{h}{6}\,
  (\mathbf{k}_1 + 2\mathbf{k}_2 + 2\mathbf{k}_3 + \mathbf{k}_4),
\end{equation}
with the usual stage definitions; its local truncation error is $O(h^5)$.
Stiff problems are routed to implicit methods (Radau, BDF)~\cite{ascher1998}.
Symplectic velocity-Verlet integration is provided for conservative mechanical
systems where long-time energy behaviour matters.

\section{Simulation Design}
The specification (\texttt{PhenomenonSpec}) records the state variables,
parameters, derivative expressions, initial conditions, integration settings,
plots, and an animation description. Derivative expressions are evaluated in a
restricted namespace exposing only NumPy~\cite{harris2020numpy} and a curated set
of mathematical functions, preventing arbitrary code execution. The code
generator renders \emph{standalone} Python (depending only on NumPy/SciPy) so
that generated projects are portable.

\section{Visualization Design}
Two visualization paths are produced. Static figures use
Matplotlib~\cite{matplotlib_docs} (time series, 2-D phase portraits, and 3-D
phase-space curves). Animations use \textsc{Manim}~\cite{manim_docs}: the
trajectory is downsampled, centred, and scaled to the frame, then drawn as a
smooth path with a moving marker; three-dimensional systems use an orbiting
camera.

\section{Results}
The framework reproduces known results quantitatively:
\begin{itemize}
  \item The simple harmonic oscillator matches $x(t) = \cos(\omega t)$ to
        better than $10^{-3}$ with RK4, and conserves energy with relative drift
        below $10^{-3}$.
  \item Exponential decay matches $e^{-t}$ to better than $10^{-4}$.
  \item The Rosenstein estimator~\cite{rosenstein1993} applied to the Lorenz
        trajectory yields a largest Lyapunov exponent of approximately $0.89$,
        close to the accepted value $\approx 0.906$, and clearly larger than the
        near-zero value for the (regular) harmonic oscillator.
\end{itemize}
A representative generated figure is included if present:
\begin{figure}[h]
\centering
\IfFileExists{../generated/lorenz-system/figure_0_phase3d.png}
  {\includegraphics[width=0.7\textwidth]{../generated/lorenz-system/figure_0_phase3d.png}}
  {\fbox{\parbox{0.7\textwidth}{\centering Generate this figure with
   \texttt{simgen generate --from-library lorenz-system --run --plot}.}}}
\caption{The Lorenz attractor produced by a generated project.}
\end{figure}

\section{Discussion}
By making the specification the single source of truth, the system guarantees
that the simulation, plots, animation, and report describe the same model. The
deterministic generator means results are reproducible: the same prompt and model
output always yield byte-identical files.

\section{Limitations}
The LLM may produce specifications that are physically plausible but
quantitatively imperfect; validation therefore relies on the offline reference
library and on automated checks. Explicit fixed-step methods can be inaccurate
for stiff systems, for which adaptive/implicit methods should be selected.
Chaotic trajectories are only predictable over a finite horizon set by the
largest Lyapunov exponent; only statistical and geometric features are robust.

\section{Future Work}
Planned extensions include native partial-differential-equation templates
(finite differences and spectral methods), automatic stiffness detection and
method selection, parameter-sweep and bifurcation tooling, and richer animation
templates (vector fields, bifurcation diagrams, and slope fields).

\section{Conclusion}
\texttt{simgen} demonstrates that natural-language prompts can be turned into
complete, reproducible scientific mini-projects by combining an LLM front end
with a deterministic, well-tested numerical and code-generation core.

\nocite{*}
\bibliographystyle{plain}
\bibliography{bibliography}

\end{document}
